\subsection{Gaussian promeasures and measures}

PROPOSITION 4. --- Let E be a locally convex space. For every positive quadratic form Q on \(E'\) , there exists one and only one promeasure \(\Gamma_{Q}\) on E such that \(F\Gamma_{Q}=e^{-Q/2}\) . The total mass of \(\Gamma_{Q}\) is equal to 1.

The uniqueness of \(\Gamma_{Q}\) follows from Prop. 3 of No.~3. The total mass of \(\Gamma_{Q}\) is equal to \((\mathcal{F}\Gamma_{\mathrm{Q}})(0)=e^{-\mathrm{Q}(0)/2}=1\) . We will prove existence in stages.

\begin{enumerate}
\def\labelenumi{\Alph{enumi})}
\tightlist
\item
  E of finite dimension n, and Q nondegenerate.
\end{enumerate}

By Lemma 2 of No.~4, the measure \(\gamma_{1}\) on R having density \(t \mapsto (2\pi)^{-1/2}e^{-t^{2}/2}\) is bounded, of total mass 1. Set \(\gamma = \gamma_{1} \otimes \cdots \otimes \gamma_{1}\) (n factors). From Lemma 2 of No.~4, one deduces

\[
\begin{array}{r l} & {\int_ {\mathbf {R} ^ {n}} e ^ {i (a _ {1} t _ {1} + \dots + a _ {n} t _ {n})} d \gamma (t _ {1}, \ldots , t _ {n}) = \prod_ {j = 1} ^ {n} \int_ {\mathbf {R}} e ^ {i a _ {j} t} d \gamma_ {1} (t)} \\ & {\qquad = \prod_ {j = 1} ^ {n} (2 \pi) ^ {- 1 / 2} \int_ {\mathbf {R}} e ^ {i a _ {j} t} e ^ {- t ^ {2} / 2} d t} \\ & {\qquad = \prod_ {j = 1} ^ {n} e ^ {- a _ {j} ^ {2} / 2}} \\ & {\qquad = \exp \big (- \frac {1}{2} (a _ {1} ^ {2} + \dots + a _ {n} ^ {2}) \big).} \end{array}
\]

Since Q is positive and nondegenerate, there exists a basis \((e_{1}^{\prime},\ldots,e_{n}^{\prime})\) of \(E^{\prime}\) orthonormal for Q (Alg., Ch. IX, §7, No.~1). Let us denote by f the isomorphism \(x\mapsto(e_{1}^{\prime}(x),\ldots,e_{n}^{\prime}(x))\) of E onto \(R^{n}\) , and by \(\Gamma_{Q}\) the measure \(f^{-1}(\gamma)\) on E. Let \(x^{\prime}=a_{1}e_{1}^{\prime}+\cdots+a_{n}e_{n}^{\prime}\) be in \(E^{\prime}\) ; then \(x^{\prime}\big(f^{-1}(t_{1},\ldots,t_{n})\big)=\sum_{j=1}^{n}t_{j}a_{j}\) for \(t_{1},\ldots,t_{n}\) real, whence

\[
\begin{array}{l} \int_ {\mathbf {E}} e ^ {i \langle x, x ^ {\prime} \rangle} d \Gamma_ {\mathbf {Q}} (x) = \int_ {\mathbf {R} ^ {n}} e ^ {i (a _ {1} t _ {1} + \dots + a _ {n} t _ {n})} d \gamma (t _ {1}, \ldots , t _ {n}) \\ \qquad = \exp \big (- \frac {1}{2} (a _ {1} ^ {2} + \dots + a _ {n} ^ {2}) = \exp \big (- \frac {1}{2} \mathbf {Q} (x ^ {\prime}) \big). \end{array}
\]

Consequently, \(\mathcal{F}\Gamma_{\mathrm{Q}} = e^{-\mathrm{Q} / 2}\) .

\begin{enumerate}
\def\labelenumi{\Alph{enumi})}
\setcounter{enumi}{1}
\tightlist
\item
  E finite-dimensional.
\end{enumerate}

Let N be the linear subspace of \(E'\) formed by the \(x'\) such that \(Q(x') = 0\) . Denote by M the orthogonal of N in E, and by j the canonical injection of M into E. The linear mapping \({}^{t}j : E' \to M'\) is surjective, with kernel N, therefore there exists on \(M'\) a nondegenerate positive quadratic form q such that \(Q = q \circ {}^{t}j\) . By the foregoing, there exists a bounded measure \(\Gamma\) on M such that \(F\Gamma = e^{-q/2}\) . Setting \(\Gamma_{Q} = j(\Gamma)\) , we have

\[
\mathcal {F} \Gamma_ {\mathrm{Q}} = (\mathcal {F} \Gamma) \circ {} ^ {t} j = \exp (- q \circ {} ^ {t} j / 2) = e ^ {- \mathrm{Q} / 2}
\]

by formula (4) of No.~3.
% semantic-fragments: legacy-input-seam
\relax{}
C) The general case.

Let \(V \in \mathcal{F}(E)\) . Denote by \(p_V\) the canonical mapping of E onto E/V, and by \(Q_V\) the positive quadratic form \(Q \circ {}^t p_V\) on \((E/V)'\) ; finally, let \(\mu_V\) be the measure on E/V with Fourier transform \(e^{-Q_V/2}\) (cf.~B)). If \(W \in \mathcal{F}(E)\) is contained in V, then \(p_V = p_{\text{VW}} \circ p_W\) , whence \(Q_V = Q_W \circ {}^t p_{\text{VW}}\) ; by formula (4) of No.~3, the measure \(p_{\text{VW}}(\mu_W)\) has as Fourier transform the function \((e^{-Q_W/2}) \circ {}^t p_{\text{VW}} = e^{-Q_V/2}\) , hence is equal to \(\mu_V\) . The family \((\mu_V)_{V \in \mathcal{F}(E)}\) is therefore a promeasure \(\mu\) on E. Formula (5) of No.~3 shows that \(\mathcal{F}\mu\) is equal to \(e^{-Q/2}\) .

DEFINITION 2. --- Let E be a locally convex space. For every positive quadratic form Q on \(E'\) , the promeasure on E whose Fourier transform is equal to \(e^{-Q/2}\) is called the Gaussian promeasure on E with variance Q, and is denoted \(\Gamma_{Q}\) . A promeasure \(\mu\) on E is said to be Gaussian if there exists a positive quadratic form Q on \(E'\) such that \(\mu = \Gamma_{Q}\) .

By an abuse of language, a bounded measure \(\mu\) on E will be said to be Gaussian with variance Q if the associated promeasure \(\widetilde{\mu}\) is equal to \(\Gamma_{Q}\) .

Remarks. --- 1) Let E be a finite-dimensional vector space, and let \(\mu\) be a positive measure on E of mass 1, such that every linear form on E belongs to \(\mathcal{L}^{2}(\mathrm{E},\mu)\) . One defines an element \(m\) of E and a positive quadratic form V on \(\mathrm{E}'\) by the formulas

\[
\langle m, x ^ {\prime} \rangle = \int_ {\mathbb {E}} \langle x, x ^ {\prime} \rangle d \mu (x), \quad \mathrm{V} (x ^ {\prime}) = \int_ {\mathbb {E}} \langle x - m, x ^ {\prime} \rangle^ {2} d \mu (x).
\]

In the traditional terminology of Probability Theory, m is called the mean and V the variance of \(\mu\) ; \(\mu\) is said to be centered if m = 0.

Now let \(a\) be an element of \(\mathbf{E}\) and \(\mathbf{Q}\) a positive quadratic form on \(\mathbf{E}'\) . Let us denote by \(\Gamma_{a,\mathbf{Q}}\) the image of the measure \(\Gamma_{\mathbf{Q}}\) under the translation \(x \mapsto x + a\) . It is easily seen that \(\Gamma_{a,\mathbf{Q}}\) is a positive measure on \(\mathbf{E}\) of mass 1, with Fourier transform \(x' \mapsto e^{i\langle a,x'\rangle - \frac{1}{2}\mathbf{Q}(x')}\) and mean \(a\) . Moreover, Prop. 6 below implies that \(\mathbf{Q}\) is the variance of \(\Gamma_{a,\mathbf{Q}}\) . One traditionally says that \(\Gamma_{a,\mathbf{Q}}\) is the Gaussian measure with mean \(a\) and variance \(\mathbf{Q}\) , and that \(\Gamma_{\mathbf{Q}} = \Gamma_{0,\mathbf{Q}}\) is the centered Gaussian measure with variance \(\mathbf{Q}\) . Since we shall only be considering centered Gaussian measures, we shall omit this qualifier.

\begin{enumerate}
\def\labelenumi{\arabic{enumi})}
\setcounter{enumi}{1}
\item
  Let \(\mathbf{Q}\) be a quadratic form on the dual \(\mathbf{E}'\) of a locally convex space \(\mathbf{E}\) . If there exists a promeasure on \(\mathbf{E}\) with Fourier transform \(e^{-\mathbf{Q} / 2}\) , the quadratic form \(\mathbf{Q}\) is necessarily positive: for, the function \(e^{-\mathbf{Q} / 2}\) is bounded on \(\mathbf{E}'\) ; therefore, for every \(x' \in \mathbf{E}'\) , the function \(t \mapsto e^{-t^2\mathbf{Q}(x') / 2} = e^{-\mathbf{Q}(tx') / 2}\) on \(\mathbf{R}\) is bounded, whence \(\mathbf{Q}(x') \geqslant 0\) .
\item
  The dual of \(\mathbf{R}\) is canonically isomorphic to \(\mathbf{R}\) and the positive quadratic forms on \(\mathbf{R}\) are the functions of the form \(t \mapsto at^2\) with \(a \geqslant 0\) . Therefore, there exists for every \(a \geqslant 0\) one and only bounded measure \(\gamma_a\) on \(\mathbf{R}\) whose Fourier transform is equal to the function \(t \mapsto e^{-at^2/2}\) ; by an abuse of language, \(\gamma_a\) is said to be the Gaussian measure on \(\mathbf{R}\) with variance \(a\) .
\end{enumerate}

The Fourier transform of \(\gamma_0\) is the constant 1, whence \(\gamma_0 = \varepsilon_0\) (unit mass at the origin of \(\mathbf{R}\) ). Suppose \(a > 0\) and denote by \(u_a\) the linear mapping \(x \mapsto a^{1/2}x\) ; then \(\mathcal{F}\gamma_a = \mathcal{F}\gamma_1 \circ {}^t u_a\) , whence \(\gamma_a = u_a(\gamma_1)\) . Lemma 2 shows that \(\gamma_1\) is the measure with density \(x \mapsto (2\pi)^{-1/2} e^{-x^2/2}\) with respect to Lebesgue measure; from this, one easily deduces

\[
d \gamma_ {a} (x) = (2 \pi a) ^ {- 1 / 2} e ^ {- x ^ {2} / 2 a} d x.\tag{15}
\]

The image of a Gaussian promeasure under a continuous linear mapping is a Gaussian promeasure. More precisely, we have the following result:

PROPOSITION 5. --- Let E and \(E_{1}\) be two locally convex spaces, and u a continuous linear mapping of E into \(E_{1}\) . Let Q be a positive quadratic form on \(E'\) , and \(Q_{1}\) the positive quadratic form \(Q \circ {}^{t}u\) on \(E_{1}'\) . Then \(u(\Gamma_{Q}) = \Gamma_{Q_{1}}\) .

Set \(\mu = u(\Gamma_{\mathbb{Q}})\) . By formula (4) of No.~3,

\[
\mathcal {F} \mu = \left(\mathcal {F} \Gamma_ {\mathrm{Q}}\right) \circ {} ^ {t} u = e ^ {- \mathrm{Q} / 2} \circ {} ^ {t} u = e ^ {- \mathrm{Q} _ {1} / 2} = \mathcal {F} \Gamma_ {\mathrm{Q} _ {1}},
\]

whence \(\mu = \Gamma_{\mathbf{Q}_1}\) by Prop. 3 of No.~3.

COROLLARY. --- Let E be a locally convex space and Q a positive quadratic form on \(E'\) . For every \(x' \in E'\) , the image of \(\Gamma_{Q}\) under \(x'\) is the Gaussian measure on R with variance \(\mathrm{Q}(x')\) .

PROPOSITION 6. --- Let E be a locally convex space, and \(\mu\) a Gaussian measure on E, with variance Q. For every integer \(n \geqslant 0\) and every \(x' \in \mathbf{E}'\) , one has the relations

\[
\int_ {\mathrm{E}} | \langle x, x ^ {\prime} \rangle | ^ {n} d \mu (x) = \pi^ {- 1 / 2} 2 ^ {n / 2} \Gamma \Big (\frac {n + 1}{2} \Big) \mathrm{Q} (x ^ {\prime}) ^ {n / 2}\tag{16}
\]

\[
\int_ {\mathrm{E}} \langle x, x ^ {\prime} \rangle^ {2 n} d \mu (x) = \frac {(2 n) !}{2 ^ {n} n !} \mathrm{Q} (x ^ {\prime}) ^ {n}\tag{17}
\]

\[
\int_ {\mathrm{E}} \langle x, x ^ {\prime} \rangle^ {2 n + 1} d \mu (x) = 0.\tag{18}
\]

In particular,

\[
\int_ {\mathrm{E}} \langle x, x ^ {\prime} \rangle^ {2} d \mu (x) = \mathrm{Q} (x ^ {\prime}) \qquad (x ^ {\prime} \in \mathrm{E} ^ {\prime}).\tag{19}
\]

If these formulas are true for an element \(x'\) of \(E'\) , then they are true for all of its multiples \(t \cdot x'\) (with t real). We may therefore content ourselves with establishing them when \(\mathrm{Q}(x')\) is equal to 0 or 1.

\begin{enumerate}
\def\labelenumi{\alph{enumi})}
\item
  Suppose \(\mathrm{Q}(x') = 0\) . The measure \(x'(\mu)\) is equal to \(\gamma_{0} = \varepsilon_{0}\) , therefore \(x'\) is zero \(\mu\) -almost everywhere; the formulas (16) to (19) are then obvious.
\item
  Suppose \(\mathrm{Q}(x') = 1\) , whence \(x'(\mu) = \gamma_{1}\) . Then
\end{enumerate}

\[
\int_ {\mathbf {E}} | \langle x, x ^ {\prime} \rangle | ^ {n} d \mu (x) = \int_ {\mathbf {R}} | t | ^ {n} d \gamma_ {1} (t) = (2 \pi) ^ {- 1 / 2} \int_ {\mathbf {R}} | t | ^ {n} e ^ {- t ^ {2} / 2} d t
\]

and (16) follows immediately from (6) (No 4, Lemma 1). Similarly, formulas (17) and (18) follow from (7) and (8). Finally, (19) is obtained by setting \(n = 1\) in (17).

We can now prove a converse of the Cor. of Prop. 5.

PROPOSITION 7. --- Let E be a locally convex space and \(\mu\) a promeasure on E. Suppose that \(x'(\mu)\) is a Gaussian measure on R for every \(x' \in E'\) . Then \(\mu\) is a Gaussian promeasure on E.

For every \(x' \in \mathrm{E}'\) , let \(\mathbf{Q}(x')\) be the variance of the Gaussian measure \(x'(\mu)\) on \(\mathbf{R}\) . One has \(x'(\mu) = \gamma_{\mathbf{Q}(x')}\) , whence

\[
(\mathcal {F} \mu) (x ^ {\prime}) = \int_ {\mathbf {R}} e ^ {i t \cdot 1} d \gamma_ {Q (x ^ {\prime})} (t) = e ^ {- Q (x ^ {\prime}) \cdot 1 ^ {2} / 2}
\]

by the definition of \(\mathcal{F}\mu\) (No.~3, formula (2)). In other words, \(\mathcal{F}\mu = e^{-Q/2}\) , and it remains to prove that \(Q\) is a positive quadratic form on \(E'\) .

For every closed linear subspace V of E with finite codimension, denote by \(p_{V}\) the canonical mapping of E onto E/V, by \(\mu_{V}\) the measure \(p_{\mathrm{V}}(\mu)\) on E/V, and set \(Q_{V} = Q \circ {}^{t} p_{V}\) . Since \(E' = \bigcup_{V \in \mathcal{F}(E)} \operatorname{Im}({}^{t} p_{V})\) and \({}^{t} p_{V}\) is injective, it suffices to prove that \(Q_{V}\) is a positive quadratic form on \((\mathrm{E}/\mathrm{V})'\) .

Let \(u \in (\mathrm{E}/\mathrm{V})'\) and \(x' = {}^{t} p_{\mathrm{V}}(u)\) . We have

\[
u (\mu_ {\mathrm{V}}) = u \big (p _ {\mathrm{V}} (\mu) \big) = x ^ {\prime} (\mu) = \gamma_ {\mathrm{Q} (x ^ {\prime})};
\]

Prop. 6 then implies

\[
\mathrm{Q} _ {\mathrm{V}} (u) = \mathrm{Q} (x ^ {\prime}) = \int_ {\mathbf {R}} t ^ {2} d \gamma_ {\mathrm{Q} (x ^ {\prime})} (t) = \int_ {\mathrm{E} / \mathrm{V}} u (x) ^ {2} d \mu_ {\mathrm{V}} (x),
\]

thus \(Q_V\) is a positive quadratic form on \((E/V)'\) .

